% GENERATED FILE. Edit canonical lessons, then run npm run generate:books.
% canonical-source-hash: fnv1a64:e4e491c4ea836b89
% canonical-lessons: RU-C255-familiya, RU-C255-vozrast, RU-C255-nomer, RU-C255-soobshcheniye, RU-C255-razgovor

\chapter{Surname, Age, Number, Message, Conversation}
\label{ch:ru-surname-age-number-message-conversation}
\hlchaptermodality{255}

\begin{chapteropening}
\textbf{By the end of this chapter:} \emph{I can say a surname, an age, a number, a message and a conversation.}

Complete the last lesson of chapter 255: I can say a surname, an age, a number, a message and a conversation.
\end{chapteropening}

\section[\texorpdfstring{\ru{фамилия} (famíliya)}{фамилия (famíliya)}]{\ru{фамилия} (famíliya) — a surname}

\label{lesson:RU-C255-familiya}



\begin{quote}
Before the new one: say the Russian for parents, then the Russian for a nephew.
\end{quote}

\subsection*{You'll want to know: \ru{фамилия}}
\textbf{\ru{фамилия}} (\emph{famíliya}) — \textquotedblleft{}a surname\textquotedblright{}.

\begin{cousinweb}
\textbf{\ru{фамилия}} — first \textquotedblleft{}a family\textquotedblright{}, from the early 1700s; it came through Polish \emph{familia} or German \emph{Familie} from Latin \emph{familia}, \textquotedblleft{}a household\textquotedblright{} (Vasmer).
\end{cousinweb}

\begin{grammarlens}[title={What you've built}]
One more word for this chapter.
\end{grammarlens}

\subsection*{Your turn}
\begin{itemize}
\raggedright
  \item \emph{Say it:} \emph{famíliya}
  \item \emph{Say it:} \emph{famíliya}, once more
  \item \emph{Say it:} say \emph{famíliya}
  \item \emph{Recall:} say the Russian for parents, then the Russian for a nephew, then say \emph{famíliya} again
\end{itemize}

\begin{tcolorbox}[breakable,colback=teal!4,colframe=teal!35!black,title={Before you move on}]
What does \ru{фамилия} mean? (\textquotedblleft{}A surname\textquotedblright{}.) Say it once more.
\end{tcolorbox}
\section[\texorpdfstring{\ru{возраст} (vózrast)}{возраст (vózrast)}]{\ru{возраст} (vózrast) — an age}

\label{lesson:RU-C255-vozrast}



\begin{quote}
Before the new one: say the Russian for a nephew, then the Russian for a surname.
\end{quote}

\subsection*{You'll want to know: \ru{возраст}}
\textbf{\ru{возраст}} (\emph{vózrast}) — \textquotedblleft{}an age\textquotedblright{}.

\begin{cousinweb}
\textbf{\ru{возраст}} — taken from Church Slavonic, says Vasmer, in place of native \emph{\ru{взрост}}: \emph{\ru{воз}-}, \textquotedblleft{}up\textquotedblright{}, with \emph{\ru{раст}}, \textquotedblleft{}growth\textquotedblright{}.
\end{cousinweb}

\begin{grammarlens}[title={What you've built}]
One more word for this chapter.
\end{grammarlens}

\subsection*{Your turn}
\begin{itemize}
\raggedright
  \item \emph{Say it:} \emph{vózrast}
  \item \emph{Say it:} \emph{vózrast}, once more
  \item \emph{Say it:} say \emph{vózrast}
  \item \emph{Recall:} say the Russian for a nephew, then the Russian for a surname, then say \emph{vózrast} again
\end{itemize}

\begin{tcolorbox}[breakable,colback=teal!4,colframe=teal!35!black,title={Before you move on}]
What does \ru{возраст} mean? (\textquotedblleft{}An age\textquotedblright{}.) Say it once more.
\end{tcolorbox}
\section[\texorpdfstring{\ru{номер} (nómer)}{номер (nómer)}]{\ru{номер} (nómer) — a number; a hotel room}

\label{lesson:RU-C255-nomer}



\begin{quote}
Before the new one: say the Russian for a surname, then the Russian for an age.
\end{quote}

\subsection*{You'll want to know: \ru{номер}}
\textbf{\ru{номер}} (\emph{nómer}) — \textquotedblleft{}a number; a hotel room\textquotedblright{}.

\begin{cousinweb}
\textbf{\ru{номер}} — \emph{\ru{нумер}} in 1704; the form with \emph{-\ru{о}-} comes from German \emph{Nummer}, from Latin \emph{numerus}, \textquotedblleft{}a number\textquotedblright{} (Vasmer). A hotel room is a numbered room.
\end{cousinweb}

\begin{grammarlens}[title={What you've built}]
One more word for this chapter.
\end{grammarlens}

\subsection*{Your turn}
\begin{itemize}
\raggedright
  \item \emph{Say it:} \emph{nómer}
  \item \emph{Say it:} \emph{nómer}, once more
  \item \emph{Say it:} say \emph{nómer}
  \item \emph{Recall:} say the Russian for a surname, then the Russian for an age, then say \emph{nómer} again
\end{itemize}

\begin{tcolorbox}[breakable,colback=teal!4,colframe=teal!35!black,title={Before you move on}]
What does \ru{номер} mean? (\textquotedblleft{}A number; a hotel room\textquotedblright{}.) Say it once more.
\end{tcolorbox}
\section[\texorpdfstring{\ru{сообщение} (soobshchéniye)}{сообщение (soobshchéniye)}]{\ru{сообщение} (soobshchéniye) — a message}

\label{lesson:RU-C255-soobshcheniye}



\begin{quote}
Before the new one: say the Russian for an age, then the Russian for a number; a hotel room.
\end{quote}

\subsection*{You'll want to know: \ru{сообщение}}
\textbf{\ru{сообщение}} (\emph{soobshchéniye}) — \textquotedblleft{}a message\textquotedblright{}.

\begin{cousinweb}
\textbf{\ru{сообщение}} comes from \emph{\ru{сообщить}}, \textquotedblleft{}to share, to report\textquotedblright{}, from \emph{\ru{общий}}, \textquotedblleft{}common\textquotedblright{}, Proto-Slavic \emph{*obĭtjĭ}, \textquotedblleft{}what is around\textquotedblright{} (Vasmer): to make something common. \emph{\ru{общий}} itself is the Church Slavonic form, with \emph{\ru{щ}}; the native Old Russian form was \emph{\ru{обчий}} (Vasmer).
\end{cousinweb}

\begin{grammarlens}[title={What you've built}]
One more word for this chapter.
\end{grammarlens}

\subsection*{Your turn}
\begin{itemize}
\raggedright
  \item \emph{Say it:} \emph{soobshchéniye}
  \item \emph{Say it:} \emph{soobshchéniye}, once more
  \item \emph{Say it:} say \emph{soobshchéniye}
  \item \emph{Recall:} say the Russian for an age, then the Russian for a number; a hotel room, then say \emph{soobshchéniye} again
\end{itemize}

\begin{tcolorbox}[breakable,colback=teal!4,colframe=teal!35!black,title={Before you move on}]
What does \ru{сообщение} mean? (\textquotedblleft{}A message\textquotedblright{}.) Say it once more.
\end{tcolorbox}
\section[\texorpdfstring{\ru{разговор} (razgovór)}{разговор (razgovór)}]{\ru{разговор} (razgovór) — a conversation}

\label{lesson:RU-C255-razgovor}



\begin{quote}
Before the new one: say the Russian for a number; a hotel room, then the Russian for a message.
\end{quote}

\subsection*{You'll want to know: \ru{разговор}}
\textbf{\ru{разговор}} (\emph{razgovór}) — \textquotedblleft{}a conversation\textquotedblright{}.

\begin{cousinweb}
\textbf{\ru{разговор}} is \emph{\ru{раз}-}, \textquotedblleft{}about, back and forth\textquotedblright{}, with \emph{\ru{говор}}, \textquotedblleft{}talk\textquotedblright{}.
\end{cousinweb}

\begin{grammarlens}[title={What you've built}]
One more word for this chapter.
\end{grammarlens}

\subsection*{Your turn}
\begin{itemize}
\raggedright
  \item \emph{Say it:} \emph{razgovór}
  \item \emph{Say it:} \emph{razgovór}, once more
  \item \emph{Say it:} say \emph{razgovór}
  \item \emph{Recall:} say the Russian for a number; a hotel room, then the Russian for a message, then say \emph{razgovór} again
\end{itemize}

\begin{tcolorbox}[breakable,colback=teal!4,colframe=teal!35!black,title={Before you move on}]
What does \ru{разговор} mean? (\textquotedblleft{}A conversation\textquotedblright{}.) Say it once more.
\end{tcolorbox}
